\NeedsTeXFormat{LaTeX2e}
\ProvidesPackage{cjhb}[29/07/2025 Cambridge Journal of Human Behaviour Article Template]

\RequirePackage{graphicx}
\RequirePackage{geometry}
\RequirePackage{xcolor}
\RequirePackage{titlesec}
\RequirePackage{fancyhdr}
\RequirePackage{hyperref}
\RequirePackage{amsthm}
\RequirePackage{fontspec}
\RequirePackage{ragged2e}
\RequirePackage{setspace}
\RequirePackage[absolute]{textpos}
\RequirePackage{calc}
\RequirePackage{rotating}
\RequirePackage{microtype}
\RequirePackage{mfirstuc}
\RequirePackage[style=apa,backend=biber]{biblatex}
\RequirePackage{csquotes}
\RequirePackage{xparse}
\RequirePackage{etoolbox}
\RequirePackage{caption}
\RequirePackage{subcaption}
\RequirePackage{multirow}
\RequirePackage{booktabs}
\RequirePackage{multicol}
\RequirePackage{endnotes}

% Hyphenation and line breaking settings
\tolerance=1000
\emergencystretch=\maxdimen
\hyphenpenalty=50
\exhyphenpenalty=50
\doublehyphendemerits=10000
\finalhyphendemerits=5000
\lefthyphenmin=3
\righthyphenmin=3
\widowpenalty=10000
\clubpenalty=10000
\brokenpenalty=10000

% ===================== VOLUME AND ISSUE SETUP =====================
\newcommand{\journalvolume}{1}
\newcommand{\journalissue}{1}
\newcommand{\setjournalvolume}[1]{\renewcommand{\journalvolume}{#1}}
\newcommand{\setjournalissue}[1]{\renewcommand{\journalissue}{#1}}

% ===================== FONT SETUP =====================
% Body font: Proforma
\setmainfont{Proforma}[
  Path=./fonts/proforma/,
  Extension = .otf,
  UprightFont = *-Regular,
  BoldFont = *-Bold,
  ItalicFont = *-Italic,
  BoldItalicFont = *-BoldItalic
]

% Title font: Helvetica Neue
\newfontfamily\articletitlefont{HelveticaNeueMedium}[
  Path=./fonts/helvetica-neue-5/,
  Extension = .otf,
  UprightFont = *,
  ItalicFont = HelveticaNeueMediumItalic,
  BoldFont = HelveticaNeueBold,
  BoldItalicFont = HelveticaNeueBoldItalic
]

% Author font: Helvetica Neue Light
\newfontfamily\authorfont{HelveticaNeueLight}[
  Path=./fonts/helvetica-neue-5/,
  Extension = .otf,
  UprightFont = *,
  ItalicFont = HelveticaNeueLightItalic,
  BoldFont = HelveticaNeueBold
]

% ========== GEOMETRY ==========
\geometry{
  paperwidth=8.27in,
  paperheight=11.69in,
  top=2.5cm,
  bottom=2cm,
  left=1.5cm,
  right=1.5cm,
  marginparwidth=1cm,
  headheight=12pt,
  headsep=1cm
}

% ===================== DISCIPLINE HANDLING =====================
\newfontfamily\disciplinefont{PT Serif}[
  Path=./fonts/PT-Serif/,
  Extension = .ttf,
  UprightFont = *-Regular
]

\newcommand{\disciplinecolor}{black}
\newcommand{\disciplinetext}{}
\newcommand{\cjhb@disciplinelist}{}

% Storage for user-defined disciplines
\newcounter{cjhbdisciplinecount}

% Helper macro to build up the discipline text
\newcommand{\cjhb@disciplinetextbuilder}{}

\newcommand{\cjhb@disciplineitem}[2]{%
  \ifx\cjhb@disciplinetextbuilder\@empty
    \gdef\cjhb@disciplinetextbuilder{{\disciplinefont\addfontfeature{LetterSpace=27}\fontsize{14}{22}\selectfont\textcolor{#1}{#2}}}%
    \renewcommand{\disciplinecolor}{#1}%
  \else
    \g@addto@macro\cjhb@disciplinetextbuilder{\\[.7em]{\disciplinefont\addfontfeature{LetterSpace=27}\fontsize{14}{22}\selectfont\textcolor{#1}{#2}}}%
  \fi
}

% User command to define a custom discipline
\newcommand{\definediscipline}[2]{%
  % #1 = discipline name (e.g., "PSYCHOLOGY")
  % #2 = RGB values as "r,g,b" (e.g., "132,75,158")
  \stepcounter{cjhbdisciplinecount}
  \expandafter\definecolor{cjhbdisc\thecjhbdisciplinecount}{RGB}{#2}
  % Store the discipline name for later retrieval
  \expandafter\gdef\csname cjhb@discname\thecjhbdisciplinecount\endcsname{#1}
  % Set the discipline color for sections to the first discipline defined
  \ifnum\value{cjhbdisciplinecount}=1
    \definecolor{sectioncolor}{RGB}{#2}%
    \renewcommand{\disciplinecolor}{sectioncolor}%
  \fi
}

% Command to set section title colour (overrides automatic first discipline colour)
\newcommand{\setsectiontitlecolour}[1]{%
  \definecolor{sectioncolor}{RGB}{#1}
  \renewcommand{\disciplinecolor}{sectioncolor}%
}

% Build banner text
\newcommand{\cjhb@buildbanner}{%
  \gdef\disciplinetext{\cjhb@drawbannertext}%
}

% Macro to draw all discipline names with their colors for banner text
\newcommand{\cjhb@drawbannertext}{%
  \newcount\tempcount
  \tempcount=1
  \loop
    \ifnum\tempcount>\value{cjhbdisciplinecount}
    \else
      \ifnum\tempcount>1\\[.7em]\fi
      {\disciplinefont\addfontfeature{LetterSpace=27}\fontsize{14}{22}\selectfont\expandafter\color{cjhbdisc\the\tempcount}\cjhb@getdisciplinename{\the\tempcount}}%
      \advance\tempcount by 1
  \repeat
}

% % Helper command to create discipline stripes
% \newcommand{\disciplinestripes}{%
%   % Use the stored discipline count
%   \newdimen\stripe@width
%   \stripe@width=1cm
%   \ifnum\value{cjhbdisciplinecount}>0
%     \divide\stripe@width by \value{cjhbdisciplinecount}
%   \fi
  
%   \makebox[0pt][l]{%
%     \raisebox{-3cm}[0pt][0pt]{%
%       % Draw the stripes for each defined discipline
%       \ifnum\value{cjhbdisciplinecount}=0
%         {\color{\disciplinecolor}\rule{\stripe@width}{4cm}}%
%       \else
%         \cjhb@drawallstripes{\stripe@width}%
%       \fi
%     }%
%   }%
% }

% Macro to draw stripes for all defined disciplines
\newcommand{\cjhb@drawallstripes}[1]{%
  \newcount\tempcount
  \tempcount=1
  \loop
    \ifnum\tempcount>\value{cjhbdisciplinecount}
    \else
      {\expandafter\color{cjhbdisc\the\tempcount}\rule{#1}{4cm}}%
      \advance\tempcount by 1
  \repeat
}

% Function to create sideways discipline text
\newcommand{\sidewaysdisciplinetext}[1]{%
  \rotatebox{#1}{%
    \parbox{40cm}{%
      \ifnum#1=-90
        \begin{flushleft}
      \else
        \begin{flushright}
      \fi
      \disciplinefont\addfontfeature{LetterSpace=27}\fontsize{11}{13}\selectfont
      \ifnum\value{cjhbdisciplinecount}=0
        {\color{\disciplinecolor}\disciplinetext}%
      \else
        \cjhb@drawalldisciplinetext%
      \fi
      \ifnum#1=-90
        \end{flushleft}
      \else
        \end{flushright}
      \fi
    }%
  }%
}

% Macro to draw all discipline names with their colors for sideways text
\newcommand{\cjhb@drawalldisciplinetext}{%
  \newcount\tempcount
  \tempcount=1
  \loop
    \ifnum\tempcount>\value{cjhbdisciplinecount}
    \else
      \ifnum\tempcount>1\hspace{1em}\fi
      {\expandafter\color{cjhbdisc\the\tempcount}\cjhb@getdisciplinename{\the\tempcount}}%
      \advance\tempcount by 1
  \repeat
}

% Helper to get discipline name by number
\newcommand{\cjhb@getdisciplinename}[1]{%
  \csname cjhb@discname#1\endcsname
}

% Main pagestripes command
\newcommand{\pagestripes}[1]{%
  \newdimen\stripe@width
  \stripe@width=1cm
  \ifnum\value{cjhbdisciplinecount}>0
    \divide\stripe@width by \value{cjhbdisciplinecount}
  \fi
  
  \begin{textblock*}{\stripe@width}(
    \if#1l
      0cm
    \else
      \ifnum\value{cjhbdisciplinecount}=1
        \dimexpr\paperwidth-\stripe@width\relax
      \else\ifnum\value{cjhbdisciplinecount}=2
        \dimexpr\paperwidth-\stripe@width-.5cm\relax
      \else\ifnum\value{cjhbdisciplinecount}=3
        \dimexpr\paperwidth-\stripe@width-.7cm\relax
      \else
        \dimexpr\paperwidth-\stripe@width-.7cm\relax
      \fi\fi\fi
    \fi,
    0pt
  )
    \ifnum\value{cjhbdisciplinecount}=0
      {\color{\disciplinecolor}\rule{\stripe@width}{4cm}}%
    \else
      \cjhb@drawallstripes{\stripe@width}%
    \fi
  \end{textblock*}%
  
  \ifnum\value{page}>1
    \begin{textblock*}{2cm}(
      \if#1l
        -0.05cm
      \else
        \dimexpr\paperwidth-1.9cm\relax
      \fi,
      4.5cm
    )
      \newcount\rot@angle
      \if#1l
        \rot@angle=-90
      \else
        \rot@angle=90
      \fi
      \sidewaysdisciplinetext{\the\rot@angle}%
    \end{textblock*}%
  \fi
}

\ProcessOptions\relax

% ========== FIXED JOURNAL BANNER ==========
\newcommand{\journalbanner}{%
  \cjhb@buildbanner
  \vspace*{-1.5cm}%
  \noindent\hspace*{-1.5cm}%
  \begin{minipage}{\dimexpr\textwidth+3cm\relax}%
    \hspace{1.8cm}%
    \begin{minipage}[t]{0.2\textwidth}
      \vspace{-.5em}%
      \hspace{-0.5cm}%
      \includegraphics[width=1\linewidth]{images/article_upper_logo.png}
    \end{minipage}%
    \begin{minipage}[t]{0.645\textwidth}
      \raggedleft 
      \vspace{1.6em}
      {\disciplinetext}
    \end{minipage}
  \end{minipage}
  \vspace{1.5em}
}

% ========== TITLE ==========
\newcommand{\articletitle}[1]{%
  
  \begingroup
    \vspace*{2em}%
    \hyphenpenalty=10000
    \exhyphenpenalty=10000
    \fontsize{23}{27}\selectfont
    \articletitlefont
    \parbox{\textwidth}{\justifying \noindent #1}
  \endgroup
}

% ========== AUTHOR AND AFFILIATION SYSTEM ==========
% Storage
\newcounter{affn}
\newcount\corraucount
\corraucount=0

% Author and affiliation lists
\def\artauthors{}
\def\auaddress{}

% Cross-reference label handling
\def\jmkLabel#1{\@bsphack\protected@write\@auxout{}{\string\Newlabel{#1}{\@currentlabel}}\@esphack}
\def\Newlabel#1#2{\expandafter\xdef\csname X@#1\endcsname{#2}}
\def\jmkRef#1{\@ifundefined{X@#1}{0}{\csname X@#1\endcsname}}

% Email storage for corresponding authors
\def\@email{}
\def\@allemail{}
\def\@correspondingemail{}

\def\email#1{%
  \gdef\@email{#1}%
}

\newcommand{\setcorrespondingemail}[1]{%
  \gdef\@correspondingemail{#1}%
}

% Author command
\let\sep\@empty
\let\authorsep\@empty

\def\author{\@ifstar\@@corrauthor\@@author}

\newcommand{\@@author}[2][]{%
  \def\@authfrstarg{#1}%
  \g@addto@macro\artauthors{%
    \def\baselinestretch{1}%
    \authorsep{#2}\unskip\ifx\@authfrstarg\empty\else\textsuperscript{\smash{{%
      \@for\@@affmark:=#1\do{\edef\affnum{\@ifundefined{X@\@@affmark}{\@@affmark}{\jmkRef{\@@affmark}}}%
      \unskip\sep\affnum\let\sep=,}}}}\fi%
    \def\authorsep{, }%
    \global\let\sep\@empty%
  }%
}

\newcommand{\@@corrauthor}[2][]{%
  \def\@authfrstarg{#1}%
  \advance\corraucount by 1%
  \g@addto@macro\artauthors{%
    \def\baselinestretch{1}%
    \authorsep{#2}\unskip\ifx\@authfrstarg\empty\else\textsuperscript{\smash{{%
      \@for\@@affmark:=#1\do{\edef\affnum{\@ifundefined{X@\@@affmark}{\@@affmark}{\jmkRef{\@@affmark}}}%
      \unskip\sep\affnum\let\sep=,}}}}\fi%
    \textsuperscript{*}%
    \def\authorsep{, }%
    \global\let\sep\@empty%
  }%
}

% Affiliation command
\newcommand{\affil}[2][]{%
  \g@addto@macro\auaddress{%
    \stepcounter{affn}%
    \xdef\@currentlabel{\theaffn}%
    \jmkLabel{\theaffn}%
    {\authorfont\fontsize{9}{12}\selectfont\textsuperscript{#1}#2\par}%
  }%
}

% Author name components
\newcommand{\fnm}[1]{\leavevmode\hbox{#1}}
\newcommand{\sur}[1]{\unskip~\leavevmode\hbox{#1}}
\newcommand{\orcidlink}[1]{\mbox{\href{https://orcid.org/#1}{\includegraphics[height=1em]{images/orcid.png}}\hspace{0.5em}}}

% Affiliation components
\newcommand{\orgdiv}[1]{#1}
\newcommand{\orgname}[1]{#1}
\newcommand{\country}[1]{#1}

% Author block
\newcommand{\neauthorblock}{%
  \vspace{2em}
  \noindent
  \begin{minipage}[t]{\textwidth}
    \begin{minipage}[t]{0.5\textwidth}
      \raggedright
      {\authorfont\fontsize{11}{12}\selectfont \artauthors\@allemail}
    \end{minipage}%
    \hfill
    \begin{minipage}[t]{0.5\textwidth}
      \raggedleft
      \copyrightstatement{\cjhb@getfirstauthorname}
    \end{minipage}
  \end{minipage}
  
  \vspace{1.5em}
  
  \noindent
  \begin{minipage}{\textwidth}
    \raggedright
    {\authorfont\fontsize{9}{12}\selectfont \auaddress}
    \ifx\@correspondingemail\empty\else
      \vspace{0.5em}
      {\authorfont\fontsize{7}{12}\selectfont *\textit{\@correspondingemail}}
    \fi
  \end{minipage}
  \vspace{.2em}
}

% Helper to get first author surname for copyright
\newcommand{\cjhb@getfirstauthorname}{%
  \article@copyrightname
}

% Command to set author heading and copyright text
\def\article@authorheading{}
\def\article@copyrightname{}
\newcommand{\setauthorheadingandcopyright}[2][]{%
  \def\article@authorheading{#2}%
  \def\article@copyrightname{\ifx\\#1\\#2\else#1\fi}%
}

% Command to format authors for header
\newcommand{\cjhb@formatauthorsforheader}{%
  \article@authorheading
}

% Copyright statement with image
\newcommand{\copyrightstatement}[1]{%
  \vspace{-.5em}
  \hbox{%
    \includegraphics[height=1em]{images/copyright.png}%
    \hspace{.3em}
    \raisebox{.6em}[0pt][0pt]{%
      \parbox[t]{\dimexpr0.8\textwidth\relax}{%
        \raggedright
        \authorfont\fontsize{6}{6}\selectfont
        \linespread{0.8}\selectfont
        \setlength{\parskip}{0pt}%
        \setlength{\baselineskip}{6pt}%
        \textcopyright\,#1. \textit{This is an Open Access article distributed under the terms of the \href{https://creativecommons.org/licenses/by-nc/4.0/}{Creative Commons Attribution-NonCommercial 4.0 License}}.%
      }%
    }%
  }%
}

% ========== REVIEW / META ==========
\newcommand{\reviewdates}[3]{%
  \authorfont\fontsize{9}{12}\selectfont
  \textbf{Received} #1\\
  \ifx\relax#2\relax
    \global\let\@seconddate\relax
  \else
    \global\def\@seconddate{#2}%
    \textbf{Revision received} #2\\%
  \fi
  \textbf{Accepted} #3%
}

\newcommand{\keywords}[1]{\authorfont\fontsize{9}{12}\selectfont\textbf{Keywords:} #1}
\newcommand{\doi}[1]{\authorfont\fontsize{9}{12}\selectfont\textbf{DOI:} \href{#1}{#1}}

% Create environments for the metadata and abstract block
\newenvironment{metablock}{%
  \hrule height 0.5pt
  \vspace{1em}
  \noindent
  \authorfont\fontsize{9}{12}\selectfont
  \begin{tabular}[t]{@{}p{0.3\textwidth}@{\hspace{0.05\textwidth}}p{0.65\textwidth}@{}}
}{%
  \end{tabular}
  \vspace{1em}
  \hrule height 0.5pt
  \vspace{-1em}
}

% Define the abstract environment to be used within metablock
\newcommand{\@seconddate}{}
\newenvironment{articleabstract}{%
  &
  \ifx\@seconddate\relax
    \vspace{-11em}%
  \else
    \vspace{-13em}%
  \fi
  \setlength{\parskip}{0.8em}%
  \setlength{\baselineskip}{12pt}%
}{}

% ========== BODY TEXT FORMATTING ==========
% Store article metadata
\def\article@title{}
\def\article@shorttitle{}

% Setters for article metadata
\newcommand{\setarticletitle}[1]{\def\article@title{#1}}
\newcommand{\setarticleshorttitle}[1]{\def\article@shorttitle{#1}}
\newcommand{\setarticledates}[1]{\def\article@dates{#1}}
\newcommand{\setarticlekeywords}[1]{\def\article@keywords{#1}}
\newcommand{\setarticledoi}[1]{\def\article@doi{#1}}
\newcommand{\setarticleabstract}[1]{\def\article@abstract{#1}}

% Command to start the main body
\newcommand{\startarticlebody}{%
  \twocolumn[{%
    \begin{@twocolumnfalse}
      \journalbanner
      \articletitle{\article@title}
      \par\vspace{.1em}%
      \neauthorblock
      \begin{metablock}
        \article@dates\\[1em]%
        \article@keywords\\[1em]%
        \article@doi%
        \begin{articleabstract}
          \article@abstract
        \end{articleabstract}
      \end{metablock}
      \vspace{2em}%
    \end{@twocolumnfalse}
  }]%
  \fontsize{9}{11}\selectfont%
  \setlength{\columnsep}{0.5cm}%
  \setlength{\parindent}{1em}%
  \setlength{\parskip}{0pt}%
}

% Custom section counter to track first paragraphs
\newcounter{firstpara}
\newcommand{\resetfirstpara}{\setcounter{firstpara}{1}}
\newcommand{\notfirstpara}{\setcounter{firstpara}{0}}

% Environment for sections that handles first paragraph indentation
\newenvironment{sectiontext}{%
  \resetfirstpara
}{%
  \notfirstpara
}

% Redefine paragraph to handle indentation
\let\oldpar\par
\DeclareRobustCommand{\par}{%
  \ifnum\value{firstpara}=1
    \setlength{\parindent}{0pt}%
    \notfirstpara
  \else
    \setlength{\parindent}{1em}%
  \fi
  \oldpar
}

% ========== SECTION FORMATTING ==========
% Variables for customising section appearance
\newcommand{\sectionfontsize}{12pt}
\newcommand{\subsectionfontsize}{10pt}
\newcommand{\sectionspacebefore}{1.5em}
\newcommand{\sectionspaceafter}{0pt}
\newcommand{\subsectionspacebefore}{1.2em}
\newcommand{\subsectionspaceafter}{0pt}

% Commands to customise section appearance
\newcommand{\setsectionfontsize}[1]{\renewcommand{\sectionfontsize}{#1}}
\newcommand{\setsubsectionfontsize}[1]{\renewcommand{\subsectionfontsize}{#1}}
\newcommand{\setsectionspacing}[2]{%
  \renewcommand{\sectionspacebefore}{#1}%
  \renewcommand{\sectionspaceafter}{#2}%
}
\newcommand{\setsubsectionspacing}[2]{%
  \renewcommand{\subsectionspacebefore}{#1}%
  \renewcommand{\subsectionspaceafter}{#2}%
}

% Section formatting
\titleformat{\section}
  {\articletitlefont\color{\disciplinecolor}\fontsize{\sectionfontsize}{14pt}\selectfont\bfseries\uppercase}
  {}
  {0pt}
  {}
  []  

\titlespacing*{\section}
  {0pt} 
  {\sectionspacebefore} 
  {0pt}  

% Subsection formatting
\titleformat{\subsection}
  {\articletitlefont\color{black}\fontsize{\subsectionfontsize}{12pt}\selectfont}
  {}
  {0pt}
  {}
  []  

\titlespacing*{\subsection}
  {0pt}  
  {\subsectionspacebefore}  
  {0pt}  

% Subsubsection formatting
\titleformat{\subsubsection}
  {\articletitlefont\color{black}\fontsize{\subsectionfontsize}{12pt}\selectfont\itshape}
  {}
  {0pt}
  {}
  []  

\titlespacing*{\subsubsection}
  {0pt}
  {\subsectionspacebefore}
  {0pt}

% ========== FIGURE AND FLOAT FORMATTING ==========
% Define new command for notes under floats
\newcommand{\figurenote}[1]{%
  \begin{minipage}{\linewidth}
    \vspace{.3em}
    \RaggedRight
    \linespread{0.85}\selectfont
    {\authorfont\fontsize{8}{9}\selectfont\itshape\textbf{Note.} #1\par}%
  \end{minipage}%
  \vspace{-1em}%
}

% Customize caption format for all float types
\DeclareCaptionFormat{cjhbfloat}{%
  {\authorfont\fontsize{9}{11}\selectfont\textbf{#1} | }\,%
  {\authorfont\fontsize{9}{11}\selectfont\itshape #3}%
}

% Set up the base style for figures
\captionsetup[figure]{
  format=cjhbfloat,
  justification=raggedright,
  singlelinecheck=true,
  labelsep=none,
  textfont=it,
  position=top,
  textformat=simple
}

% Make figure* use exactly the same format as figure
\captionsetup[figure*]{
  format=cjhbfloat,
  justification=raggedright,
  singlelinecheck=true,
  labelsep=none,
  textfont=it,
  position=top,
  textformat=simple
}

% Apply same settings to tables
\captionsetup[table]{format=cjhbfloat}
\captionsetup[table*]{format=cjhbfloat}
\captionsetup[subfigure]{format=cjhbfloat}
\captionsetup[subtable]{format=cjhbfloat}

% Set table text to use Helvetica Light font
\renewcommand{\arraystretch}{1.2}
\renewcommand{\tabcolsep}{6pt}

% Customize table environment to use Helvetica Light
\let\oldtabular\tabular
\renewcommand{\tabular}{\authorfont\fontsize{8}{8}\oldtabular}

% Ensure equation numbers follow similar style
\renewcommand{\theequation}{%
  {\authorfont\fontsize{9}{11}\selectfont\arabic{equation}}%
}

% ========== PAGE FOOTER ==========
\pagestyle{fancy}
\fancyhf{}

% Set up the TPGrid for absolute positioning
\setlength{\TPHorizModule}{1mm}
\setlength{\TPVertModule}{1mm}
\textblockorigin{0mm}{0mm}

% Headers for all pages
\fancyhead[L]{%
  \ifnum\value{page}>1
    {\authorfont\fontsize{8}{10}\selectfont\makebox[\dimexpr0.45\textwidth][l]{\cjhb@formatauthorsforheader}}%
  \fi
  \ifodd\value{page}
    \hbox to 0pt{\pagestripes{l}\hss}%
  \fi
}
\fancyhead[R]{%
  \ifnum\value{page}>1
    {\authorfont\fontsize{8}{10}\selectfont\makebox[\dimexpr0.45\textwidth][r]{\article@shorttitle}}%
  \fi
  \ifodd\value{page}\else
    \hbox to 0pt{\pagestripes{r}\hss}%
  \fi
}

% Footers for all pages
\fancyfoot[L]{\authorfont\fontsize{7}{9}\selectfont\itshape Cambridge Journal of Human Behaviour}
\fancyfoot[C]{\authorfont\fontsize{7}{9}\selectfont\bfseries\thepage}
\fancyfoot[R]{\authorfont\fontsize{7}{9}\selectfont\itshape Volume \journalvolume\upshape\ | Issue \journalissue}

% Conditional header rule
\renewcommand{\headrule}{%
  \ifnum\value{page}>1
    \hrule width\headwidth height0.5pt \vskip-0.5pt
  \fi
}

\renewcommand{\footrulewidth}{.5pt}

% Plain style for first page
\fancypagestyle{plain}{%
  \fancyhf{}
  \fancyhead[L]{\ifodd\value{page}\hbox to 0pt{\pagestripes{l}\hss}\fi}
  \fancyhead[R]{\ifodd\value{page}\else\hbox to 0pt{\pagestripes{r}\hss}\fi}
  \fancyfoot[L]{\authorfont\fontsize{7}{9}\selectfont\itshape Cambridge Journal of Human Behaviour}
  \fancyfoot[C]{\authorfont\fontsize{7}{9}\selectfont\bfseries\thepage}
  \fancyfoot[R]{\authorfont\fontsize{7}{9}\selectfont\itshape Volume \journalvolume\upshape\ | Issue \journalissue}
  \renewcommand{\headrulewidth}{0pt}
  \renewcommand{\footrulewidth}{.5pt}
}

% ========== PARAGRAPH ==========
\setlength{\parindent}{0pt}
\setlength{\parskip}{0.8em}

% ========== QUOTE ENVIRONMENT ==========
\renewenvironment{quote}{%
  \list{}{%
    \rightmargin0pt%
    \leftmargin3em%
  }%
  \item\relax\justifying%
}{%
  \endlist
}

% ========== FOOTNOTE HANDLING ==========
% Redefine footnote to use endnotes
\let\footnote\endnote

% Remove the default endnotes title
\renewcommand{\notesname}{}

% Custom formatting for the Notes section
\newcommand{\printnotes}{%
  \ifnum\value{endnote}>0%
    \section*{Notes}%
    \vspace{-1.5em}%
    \begingroup
      \parindent 0pt%
      \parskip 0.3em%
      \theendnotes%
    \endgroup
  \fi
}

% Customise note numbers to be superscript in both text and Notes section
\renewcommand\makeenmark{\textsuperscript{\theenmark}}

% Make endnote numbers in text superscripted
\renewcommand\@makeenmark{\textsuperscript{\theenmark}}

% Remove the period after note numbers in the Notes section and any heading
\renewcommand\enoteformat{%
  \noindent%
  \theenmark\hspace{0.5em}%
}

% ========== INTERDISCIPLINARY COMMENTARIES ==========
% Commands for creating interdisciplinary commentaries
\newcommand{\commentarydiscipline}[2]{%
  \definecolor{commentarytempcolor}{RGB}{#1}%
  \vspace{1em}%
  \noindent\hfill%
  {\disciplinefont\fontsize{10}{12}\selectfont\bfseries\addfontfeature{LetterSpace=27}\color{commentarytempcolor}#2}%
  \par\noindent\hfill%
  {\color{commentarytempcolor}\rule{0.4\textwidth}{3pt}}%
  \par\vspace{2em}%
}

\newcommand{\commentarytitle}[1]{%
  \noindent%
  {\articletitlefont\fontsize{12}{16}\selectfont\bfseries #1}%
  \par\vspace{1.5em}%
}

\newcommand{\commentaryauthor}[1]{%
  \noindent%
  {\authorfont\fontsize{10}{11}\selectfont #1}%
  \par\vspace{0.1em}%
}

\newcommand{\commentaryaffiliation}[1]{%
  \noindent%
  {\authorfont\fontsize{8}{10}\selectfont #1}%
  \par\vspace{1.5em}%
  \noindent\hrule height 1pt%
  \par\vspace{1em}%
}

\newcommand{\commentaryabstract}[1]{%
  \noindent%
  {\authorfont\fontsize{9}{11}\selectfont #1}%
  \par\vspace{1em}%
  \noindent\hrule height 1pt%
  \par\vspace{1em}%
}

% Environment for commentary content that uses subsection formatting
\newenvironment{commentarycontent}{%
  \begin{sectiontext}%
}{%
  \end{sectiontext}%
  \vspace{2em}%
}

% Command to start the commentaries section
\newcommand{\startcommentaries}{%
  \vspace{4em}%
  \noindent\hrule height 3pt width\columnwidth%
  \vspace{2em}%
  \noindent{\articletitlefont\fontsize{17}{18}\selectfont Interdisciplinary Commentaries}%
  \par\vspace{2em}%
}

% ========== BIBLIOGRAPHY FORMATTING ==========
\renewcommand*{\bibfont}{\authorfont\fontsize{7}{9}\selectfont}

\setlength{\bibhang}{0.5em}
\setlength{\bibitemsep}{0.2em}

\DefineBibliographyStrings{english}{%
  references = {Article references}%
}

\DeclareSourcemap{
  \maps[datatype=bibtex]{
    \map{
      \perdatasource{references/references.bib}
      \step[fieldset=keywords, fieldvalue={cjhbarticleref}, append]
    }
    \map{
      \perdatasource{references/commentaryreferences.bib}
      \step[fieldset=keywords, fieldvalue={cjhbcommentaryref}, append]
    }
  }
}

% Command to print article references
\newcommand{\printarticlereferences}{%
  \clearpage
  \onecolumn
  \noindent{\authorfont\color{black}\fontsize{20pt}{14pt}\selectfont Article references}\par
  \vspace{.5em}%
  \begin{multicols}{3}
  \printbibliography[heading=none, keyword=cjhbarticleref]
  \end{multicols}%
}

% Command to load commentary references
\newcommand{\loadcommentaryreferences}{%
  \addbibresource{references/commentaryreferences.bib}%
}

% Command to print commentary references
\newcommand{\printcommentaryreferences}{%
  \vspace{2em}%
  \noindent\hrule height 1pt width\textwidth%
  \vspace{2em}%
  \noindent{\authorfont\color{black}\fontsize{20pt}{14pt}\selectfont Commentary references}\par
  \vspace{.5em}%
  \begin{multicols}{3}
  \printbibliography[heading=none, keyword=cjhbcommentaryref]
  \end{multicols}%
}
